\section{Discussion}
\markright{Discussion}

The forward-bias readings show the nonlinear behavior of the silicon p-n junction diode.
At low applied voltage, the current is very small because the external voltage is not
large enough to overcome the potential barrier of the depletion region. With the
$10\,\text{k}\Omega$ series resistor used in the experiment, the current remained in the
low-current region, rising from about $0.001\,\text{mA}$ to $0.102\,\text{mA}$ as the
diode voltage increased from $0.29\,\text{V}$ to $0.48\,\text{V}$. At higher current
levels, a silicon diode normally approaches the familiar forward drop of about
$0.7\,\text{V}$.

In reverse bias, the measured current remains very small even when the reverse voltage is
increased. This happens because the depletion region widens and blocks majority carrier
flow. Only a small reverse saturation current flows due to minority carriers. Since the
applied reverse voltage was kept below breakdown, no sudden increase in reverse current
was observed.

The simulation results agree with the practical characteristic. Minor differences between
simulated and observed values may occur because a real diode has internal resistance,
temperature dependence, contact resistance, meter loading, and manufacturing tolerance.
The value of the series resistor strongly affects the current. Since a
$10\,\text{k}\Omega$ resistor was used, the current stayed small even when the applied
voltage was increased.

Overall, the experiment verifies that a silicon p-n junction diode conducts effectively
only in forward bias after the cut-in voltage and behaves almost as an open circuit in
reverse bias under normal operating conditions.

